\chapter{Background and Literature Review}\label{ch:background}

\section{Theoretical Foundations}
Present the theory with numbered equations that you refer to by label:
\begin{equation}\label{eq:energy}
  E = mc^{2}.
\end{equation}
\Cref{eq:energy} is numbered by chapter. Use \texttt{\textbackslash vect}, \texttt{\textbackslash mat} and
\texttt{\textbackslash qty} for notation and units, for example $\vect{F}=m\vect{a}$ and \qty{9.81}{\metre\per\second\squared}.

\begin{definitionbox}[Definition]
A boxed definition for the central concept of the chapter.
\end{definitionbox}

\begin{theorem}[Pythagoras]\label{thm:pyth}
For a right triangle with legs $a,b$ and hypotenuse $c$, $a^{2}+b^{2}=c^{2}$.
\end{theorem}
\begin{proof}
Standard; see \cite{knuth1984}.
\end{proof}

\section{Related Work}
Compare the main existing approaches in a table.
\begin{center}\atlrows
  \captionof{table}{Comparison of existing approaches.}\label{tab:related}
  \begin{tabular}{@{}llll@{}}
    \rowcolor{atl-primary}\atlhead{Approach} & \atlhead{Strength} & \atlhead{Weakness} & \atlhead{Ref.}\\
    Method A & Fast & Inaccurate & \cite{knuth1984} \\
    Method B & Accurate & Slow & \cite{lamport1994} \\
  \end{tabular}
\end{center}

\section{Chapter Summary}
Summarise what the reader should take from this chapter.
